\section*{Bloque A. Fundamentos teóricos y conceptuales}

\noindent\textit{Preguntas 1 a 30. Evalúan la comprensión analítica de externalidades en producción y consumo, costos sociales vs privados, impuestos pigouvianos, Teorema de Coase, bienes públicos puros, condición de Samuelson, equilibrio de Lindahl, recursos comunes y modelos de selección adversa y riesgo moral.}

\q{En la teoría microeconómica del bienestar, una externalidad se define formalmente como:}
\begin{opts}
    \item fluctuaciones exógenas en la tasa de interés real alteran los precios relativos de mercado.
    \item un monopolio natural fija una tarifa regulatoria por encima de su costo marginal de corto plazo.
    \item la acción de un agente afecta directamente el bienestar o producción de otro sin mediar precio.
    \item la aplicación de un impuesto indirecto distorsiona las decisiones de consumo de los hogares.
\end{opts}

\q{En presencia de una externalidad negativa en la producción (por ejemplo, vertimiento de efluentes contaminantes), la relación entre el costo marginal social ($CMgS$) y el costo marginal privado ($CMgP$) es:}
\begin{opts}
    \item $CMgS = CMgP + CMgE$, superando el costo social al costo privado.
    \item $CMgS = 0$ exactamente en el nivel de producción óptimo paretiano.
    \item $CMgS = CMgP$, al carecer de valoración monetaria el entorno natural.
    \item $CMgS = CMgP - CMgE$, descontando el daño ambiental marginado.
\end{opts}

\q{En un mercado competitivo libre sin regulación donde existe una externalidad negativa en la producción, el equilibrio de mercado genera:}
\begin{opts}
    \item producción superior al óptimo social ($Q_C > Q^*$), generando exceso.
    \item un precio de mercado estrictamente superior al costo marginal social total.
    \item una asignación eficiente idéntica a la que maximiza el excedente neto social.
    \item un volumen de producción inferior al óptimo social, causando escasez crónica.
\end{opts}

\q{Cuando existe una externalidad positiva en el consumo (por ejemplo, vacunación o educación sanitaria sobre agua potable):}
\begin{opts}
    \item la política óptima exige aplicar un impuesto específico a los consumidores.
    \item el beneficio marginal social supera al privado ($BMgS > BMgP$): subprovisión.
    \item el mercado competitivo genera una cantidad excesiva frente al óptimo paretiano.
    \item el costo marginal social supera ampliamente al costo marginal privado.
\end{opts}

\q{Un impuesto pigouviano sobre la producción de un bien que genera externalidad negativa debe fijarse exactamente igual a:}
\begin{opts}
    \item el costo marginal externo evaluado en el nivel socialmente óptimo ($t^* = CMgE(Q^*)$).
    \item el presupuesto total de remediación ambiental dividido por el volumen producido.
    \item el excedente neto del consumidor calculado en el equilibrio libre sin intervención.
    \item la diferencia absoluta entre el precio competitivo y el costo medio mínimo de planta.
\end{opts}

\q{Para corregir una externalidad positiva, la intervención pigouviana estándar consiste en:}
\begin{opts}
    \item otorgar subsidio unitario equivalente al beneficio externo ($s^* = BMgE(Q^*)$).
    \item establecer un precio máximo vinculante inferior al costo marginal privado.
    \item imponer una cuota máxima de producción para evitar saturar la demanda de mercado.
    \item prohibir la participación privada en el sector y estatizar íntegramente la oferta.
\end{opts}

\q{El Teorema de Coase establece que en presencia de externalidades, la negociación voluntaria entre las partes privadas conduce a una asignación eficiente de recursos si se cumplen los supuestos de:}
\begin{opts}
    \item intervención directa del regulador mediante tarifas máximas y subsidios cruzados.
    \item existencia de un monopolio natural regulado con rendimientos crecientes a escala.
    \item derechos de propiedad bien definidos, sin costos de transacción e información plena.
    \item bienes perfectamente indivisibles y ausencia total de bienes sustitutos cercanos.
\end{opts}

\q{Según el Teorema de Coase, la asignación final de recursos y el nivel de contaminación de equilibrio:}
\begin{opts}
    \item dependen críticamente de a cuál de las partes se asigne el derecho de propiedad.
    \item serán siempre ineficientes si el derecho inicial recae en los agentes afectados.
    \item convergen invariablemente hacia un nivel de contaminación ambiental igual a cero.
    \item son independientes de quien reciba el derecho inicial, cambiando solo la renta.
\end{opts}

\q{La principal limitación práctica para la aplicación del Teorema de Coase en problemas ambientales a gran escala (como la contaminación hídrica de una cuenca hidrográfica) radica en:}
\begin{opts}
    \item la ausencia de series históricas estadísticas para estimar elasticidades de demanda.
    \item la dificultad contable de valorar los gastos fijos operativos de las empresas.
    \item altos costos de transacción, asimetría informativa y free-riding masivo.
    \item la imposibilidad jurídica absoluta de delimitar derechos de aprovechamiento hídrico.
\end{opts}

\q{En un sistema de derechos de emisión transferibles (Cap-and-Trade) para controlar la contaminación:}
\begin{opts}
    \item el regulador fija el tope agregado de emisiones y el mercado determina el precio del permiso.
    \item los derechos de emisión se reasignan diariamente mediante subasta gratuita no vinculante.
    \item el regulador fija el precio unitario del permiso y el mercado ajusta la contaminación total.
    \item todas las empresas reducen emisiones en idéntico porcentaje sin posibilidad de intercambio.
\end{opts}

\q{Bajo criterios de eficiencia en costos (cost-effectiveness), los instrumentos basados en el mercado (impuestos y permisos) son preferibles a la regulación tradicional de 'mando y control' (estándares uniformes) porque:}
\begin{opts}
    \item eliminan la necesidad técnica de monitorear y fiscalizar las descargas efectivas en campo.
    \item garantizan una recaudación fiscal proporcional al tamaño del producto bruto sectorial.
    \item igualan los costos marginales de abatimiento entre firmas, minimizando el costo agregado.
    \item obligan a todas las empresas contaminantes a reducir idéntica proporción sin excepciones.
\end{opts}

\q{Un Bien Público Puro se caracteriza rigurosamente por cumplir dos propiedades fundamentales:}
\begin{opts}
    \item rivalidad estricta en el consumo y libre acceso al recurso.
    \item no rivalidad en el consumo pero facilidad técnica de exclusión por precio.
    \item producción exclusiva a cargo de empresas con propiedad estatal.
    \item no rivalidad en el consumo y no exclusión en el acceso.
\end{opts}

\q{En la tipología estándar de bienes según rivalidad y exclusión, un 'Recurso Común' (o bien comunal) es:}
\begin{opts}
    \item rival en el consumo y no excluible en su acceso.
    \item rival en consumo y fácilmente excluible por precio.
    \item no rival y no excluible de forma simultánea.
    \item no rival en consumo pero técnicamente excluible.
\end{opts}

\q{El fenómeno del 'polizón' (Free-Rider Problem) en la provisión de bienes públicos ocurre principalmente debido a:}
\begin{opts}
    \item la imposibilidad de excluir a quien no paga, incentivando ocultar su valoración real.
    \item la existencia de barreras legales de entrada asociadas a monopolios naturales.
    \item la carencia técnica de sistemas de medición y facturación electrónica individualizada.
    \item la pendiente positiva de la curva de oferta agregada bajo competencia perfecta.
\end{opts}

\q{La condición de eficiencia de Samuelson para la provisión óptima de un bien público puro ($G$) frente a un bien privado ($x$) establece que:}
\begin{opts}
    \item el costo marginal de proveer una unidad adicional del bien público sea exactamente nulo.
    \item la suma de las $RMS$ individuales iguale a la transformación ($\sum RMS_i = RMT$).
    \item la cuota uniforme pagada por cada ciudadano iguale al costo medio total de provisión.
    \item cada relación marginal de sustitución individual iguale la de transformación ($RMS_i = RMT$).
\end{opts}

\q{La construcción geométrica de la curva de demanda de mercado de un bien público difiere de la de un bien privado en que:}
\begin{opts}
    \item presenta pendiente positiva en todo su tramo por efecto riqueza y paradoja de Giffen.
    \item se obtiene sumando verticalmente las disposiciones marginales a pagar por cada cantidad.
    \item se construye sumando horizontalmente las cantidades individuales demandadas a cada precio.
    \item resulta perfectamente elástica al nivel exacto del costo marginal constante de provisión.
\end{opts}

\q{En el concepto del Equilibrio de Lindahl para la financiación de un bien público:}
\begin{opts}
    \item todos los residentes tributan una tarifa fija uniforme determinada por ley orgánica.
    \item la provisión del bien depende únicamente de contribuciones voluntarias y anónimas.
    \item el gasto público se financia exclusivamente mediante impuestos sobre rentas monopólicas.
    \item cada agente paga un precio personalizado igual a su beneficio marginal en el óptimo.
\end{opts}

\q{El mecanismo de incentivos de Vickrey-Clarke-Groves (VCG) resuelve el problema de revelación de preferencias en bienes públicos mediante:}
\begin{opts}
    \item la prohibición normativa de revelar valoraciones monetarias en consultas vecinales.
    \item un impuesto pivote (Clarke) igual al daño externo que su voto impone sobre los demás.
    \item el sorteo aleatorio de la carga tributaria entre los votantes registrados del distrito.
    \item la privatización directa de toda infraestructura comunal mediante licitación pública.
\end{opts}

\q{La 'Tragedia de los Comunes' descrita por Garrett Hardin (1968) en el uso de recursos de libre acceso se explica microeconómicamente porque:}
\begin{opts}
    \item los usuarios perciben el recurso natural como un bien inferior de consumo decreciente.
    \item la autoridad sectorial prohíbe adoptar tecnologías eficientes de extracción sostenible.
    \item los costos marginales privados de extracción tienden a infinito en todo nivel de uso.
    \item cada agente extrae hasta igualar beneficio marginal y costo medio, ignorando congestión.
\end{opts}

\q{Un 'Bien Meritorio' (Merit Good), concepto introducido por Richard Musgrave, se distingue porque:}
\begin{opts}
    \item el Estado lo promueve o subsidia pues los agentes subestiman sus beneficios sociales.
    \item el mercado privado competitivo lo asigna óptimamente sin requerir subsidios estatales.
    \item es demandado exclusivamente por hogares pertenecientes a los deciles más altos de ingreso.
    \item exhibe costos marginales rápidamente crecientes y elasticidad ingreso negativa.
\end{opts}

\q{En la teoría de la información asimétrica, la Selección Adversa se diferencia del Riesgo Moral en que:}
\begin{opts}
    \item la primera aborda acciones ocultas ex-post; el segundo, características ocultas ex-ante.
    \item ambos conceptos describen fallas idénticas de colusión horizontal entre aseguradoras.
    \item el riesgo moral ocurre únicamente en firmas privadas y la selección en consumidores.
    \item la primera trata características ocultas ex-ante; el segundo, acciones ocultas ex-post.
\end{opts}

\q{En el célebre modelo del 'Mercado de Cacharros' (The Market for Lemons) de George Akerlof (1970):}
\begin{opts}
    \item al ofrecerse un precio promedio por calidad desconocida, se expulsan los autos buenos.
    \item el precio de equilibrio converge siempre al valor de los vehículos de mayor calidad.
    \item las unidades de alta calidad desplazan por competencia a los vehículos defectuosos.
    \item la asimetría informativa desaparece sin costo mediante garantías verbales de venta.
\end{opts}

\q{En el modelo de Selección Adversa en seguros de salud y vida:}
\begin{opts}
    \item los copagos quedan prohibidos por generar pérdidas irrecuperables de eficiencia.
    \item los clientes de menor riesgo intrínseco demandan la máxima cobertura disponible.
    \item quienes tienen mayor riesgo compran más seguros, subiendo primas y expulsando sanos.
    \item las aseguradoras obtienen ganancias ilimitadas aun fijando primas uniformes sin deducible.
\end{opts}

\q{En el modelo de Señalización de Michael Spence (1973), para que un título educativo funcione como una señal informativa creíble de alta productividad se requiere que:}
\begin{opts}
    \item educarse cueste menos a los trabajadores más productivos (Single-Crossing).
    \item la educación eleve obligatoriamente la productividad física en idéntica proporción.
    \item el Estado financie de manera universal e irrestricta toda la educación superior.
    \item el costo marginal educativo sea exactamente idéntico para todos los trabajadores.
\end{opts}

\q{En el modelo de Cribado (Screening) de Rothschild y Stiglitz (1976) en mercados de seguros, la empresa aseguradora desinformada induce la autorrevelación de los clientes ofreciendo:}
\begin{opts}
    \item un menú separador: cobertura parcial barata para bajo riesgo y total cara para alto riesgo.
    \item un contrato agrupador único con cobertura completa al costo actuarial medio poblacional.
    \item una póliza uniforme asignada mediante subasta descendente abierta a todo postor.
    \item una cobertura básica gratuita financiada íntegramente con impuestos sobre la renta.
\end{opts}

\q{En un modelo de contratos con asimetría de información, un 'Equilibrio Agrupador' (Pooling Equilibrium) se caracteriza porque:}
\begin{opts}
    \item cada tipo de agente elige un contrato distinto ajustado a su perfil de riesgo privado.
    \item todas las compañías participantes operan en equilibrio con pérdidas actuariales netas.
    \item todos los tipos eligen idéntico contrato o acción, ocultando su información privada.
    \item ningún individuo con aversión al riesgo decide participar en el mercado de pólizas.
\end{opts}

\q{El problema del Riesgo Moral (Moral Hazard) en las relaciones laborales entre un empleador (Principal) y un trabajador (Agente) surge cuando:}
\begin{opts}
    \item el salario contractual iguala exactamente el producto marginal del trabajo observable.
    \item ambas partes exhiben funciones de utilidad estrictamente neutrales ante el riesgo.
    \item el esfuerzo del agente no es observable ni verificable y el resultado depende del azar.
    \item el agente posee mayor dotación inicial de capital financiero que el propio principal.
\end{opts}

\q{En la teoría clásica de incentivos bajo aversión al riesgo, el trade-off fundamental del contrato óptimo Principal-Agente se da entre:}
\begin{opts}
    \item discriminación intertemporal de precios y arbitraje geográfico espacial.
    \item economías de escala productivas y barreras tecnológicas a la entrada.
    \item recaudación tributaria óptima y cumplimiento de metas ambientales locales.
    \item incentivar el esfuerzo y asignar riesgo al actor más tolerante.
\end{opts}

\q{En los contratos de seguros de propiedad o automóviles, la inclusión de deducibles y copagos sirve primordialmente para:}
\begin{opts}
    \item evitar que el beneficiario solicite peritajes independientes de valoración de daños.
    \item suprimir totalmente los costos administrativos y de monitoreo de la aseguradora.
    \item mitigar el riesgo moral haciendo asumir al asegurado parte del siniestro para prevenir.
    \item garantizar que la prima actuarial neta converja a cero en periodos sin accidentes.
\end{opts}

\q{La teoría de los Salarios de Eficiencia de Shapiro y Stiglitz (1984) demuestra que pagar un salario superior al de equilibrio de mercado:}
\begin{opts}
    \item elimina todo desempleo involuntario haciendo que el mercado de trabajo se vacíe.
    \item disciplina al trabajador frente a la holgazanería al elevar el costo real del despido.
    \item reduce el producto medio laboral por generar relajamiento en la supervisión interna.
    \item elimina por completo la selección adversa sobre la calidad de insumos intermedios.
\end{opts}

